%\documentclass[12pt,oneside]{amsart}
\documentclass{scrartcl}

%\documentclass{article}

\usepackage{amsthm,amsmath,amssymb}
\usepackage[numbers]{natbib}

%\usepackage[colorlinks]{hyperref}
\usepackage{hyperref}
\hypersetup{
    colorlinks,%
    citecolor=black,%
    filecolor=black,%
    linkcolor=black,%
    urlcolor=black
}
\usepackage{hypernat}
\usepackage[top=2.5cm, bottom=2.5cm, left=2.8cm,
right=2.8cm]{geometry}
\usepackage{graphicx}


%\setlength{\parskip}{0pt}
%\setlength{\parsep}{0pt}
%\setlength{\headsep}{0pt}
%\setlength{\topskip}{0pt}
%\setlength{\topmargin}{0pt}
%\setlength{\topsep}{0pt}
%\setlength{\partopsep}{0pt}

\linespread{1}

%\usepackage{marvosym}

%\textwidth = 6.5 in \textheight = 9.2 in \oddsidemargin = 0.0 in
%\evensidemargin = 0.0 in \topmargin = -0.0 in \headheight = 0.0 in
%\headsep = 0.0in \parskip = 0.0in \parindent = 0.0in
%\def\baselinestretch{0.871}


\newtheorem{theorem}{Theorem}[section]
\newtheorem{proposition}[theorem]{Proposition}
\newtheorem{problem}[theorem]{Problem}
\newtheorem{fact}[theorem]{Fact}
\newtheorem{lemma}[theorem]{Lemma}
\newtheorem{defn}[theorem]{Definition}
\newtheorem{corollary}[theorem]{Corollary}
\newtheorem{remark}{Remark}[section]
\newtheorem{example}[theorem]{Example}

\newcommand{\E}{{\mathbb E}}
\newcommand{\se}{{\mathcal{E}}}
\newcommand{\R}{{\mathbb R}}
\renewcommand{\P}{{\mathbb P}}
\newcommand{\ac}{{\mathcal{A}}}
\newcommand{\A}{{\mathcal{A}}}
\newcommand{\B}{{\mathcal{B}}}
\newcommand{\C}{{\mathcal{C}}}
\newcommand{\M}{{\mathcal{M}}}
\newcommand{\Mf}{{\mathfrak{M}}}
\newcommand{\T}{{\mathcal{T}}}
\newcommand{\Z}{{\mathcal{Z}}}
\newcommand{\K}{{\mathcal{K}}}
\newcommand{\lc}{{\mathcal{L}}}
\newcommand{\Ps}{{\mathcal{P}}}
\newcommand{\G}{{\mathcal{G}}}
\newcommand{\pa}{{\mathring{p}}}
\newcommand{\F}{{\cal F}}
\newcommand{\samp}{{\mathcal{X}}}
\newcommand{\X}{{\mathcal{X}}}
\newcommand{\hi}{{\hat{\Phi}}}
\newcommand{\data}{{\text{data}}}
\newcommand{\relu}{{\text{ReLU}}}

\newcommand{\ridge}{{\hat{\beta}_{\text{ridge}}(\lambda)}}
\newcommand{\lasso}{{\hat{\beta}_{\text{lasso}}(\lambda)}}
\newcommand{\ridgemu}{{\hat{\mu}_t^{\text{ridge}}(\lambda)}}
\newcommand{\lassomu}{{\hat{\mu}_t^{\text{lasso}}(\lambda)}}


\newcounter{rcnt}[section]
\renewcommand{\thercnt}{(\roman{rcnt})}

\def\qt#1{\qquad\text{#1}}

\def\argmin{\mathop{\rm argmin}}
\def\argmax{\mathop{\rm argmax}}


\setlength{\parskip}{1.4 \medskipamount}
%\sloppy
%\linespread{1.3}

\begin{document}

\title{STAT 238 - Bayesian Statistics  \\
  Lecture Twenty} 
\subtitle{Spring 2026, UC Berkeley}

\author{Aditya Guntuboyina}


\date{11 March 2026}

\maketitle

\section{Bayesian Inference in a High-Dimensional Linear Regression
  Model}

We studied the following model in the last lecture.

We have a response variable $y$ and a single covariate $x$. In our
example, $y$ denotes weekly earnings and $x$ denotes years of
experience. The covariate $x$ takes the values $0, 1, \dots, m$ for
some fixed integer $m$. 

Our data is $(x_i, y_i), i = 1, \dots, n$. The model is:
\begin{align}\label{ibm_mod}
  y_i = \beta_0 + \beta_1 x_i + \beta_2 \relu(x_i - 1) + \dots +
  \beta_{m}\relu(x_i - (m-1)) + \epsilon_i
\end{align}
where $\relu(u) := \max(u, 0)$, and $\epsilon_i
\overset{\text{i.i.d}}{\sim} N(0, \sigma^2)$. Note we did not include
$\relu(x_i - m)$ because it always equals 0.

The model \eqref{ibm_mod} can be rewritten in the usual regression way
as:
\begin{align*}
  y = X \beta + \epsilon ~~\text{with $\epsilon \sim N(0, \sigma^2
  I_d)$}. 
\end{align*}
Here $X$ is the $n \times (m+1)$ matrix with columns $1, x,
\relu(x-j)$ for $j = 1, \dots, m-1$, where $x$ denotes observed values
of the experience variable. 

The usual least squares analysis coincides with Bayesian inference
using the prior:
\begin{align}\label{lsp}
  \beta_0, \dots, \beta_m \overset{\text{i.i.d}}{\sim} N(0, C)
  \text{ and } \log \sigma \sim \text{uniform}(-C, C). 
\end{align}
for $C \rightarrow \infty$. In Problem 7 of Homework 3, it is proved
that the posterior for the above prior is: 
\begin{align*}
  \beta \mid \text{data}, \sigma \sim N_{m+1}(\hat{\beta}, \sigma^2
  (X^T X)^{-1}) ~ \& ~  f_{\sigma \mid \text{data}}(\sigma)
  \propto \sigma^{-n+m} \exp \left(-\frac{y^T y - y^T X (X^T X)^{-1}
  X^T y}{2
  \sigma^2} \right) I\{\sigma > 0\}
\end{align*}
where $\hat{\beta} = (X^T X)^{-1} X^T y$ is the least squares
estimator. If we marginalize $\sigma$ and write the posterior
distribution of $\beta$, we will get a $t$-distribution. The posterior
distribution for $\sigma$ above can be written in terms of the Gamma
(or chi-squared) distribution (this is problem 7(d) in Homework 3) as:
\begin{align*}
  \frac{1}{\sigma^2} \Big| \text{data} \sim
 \text{Gamma} \left(\frac{n-m-1}{2}, \frac{y^T y -
  y^T X (X^T X)^{-1} X^T y}{2}
  \right). 
\end{align*}
Because least squares does not give
sensible results in this example, we change the prior in \eqref{lsp}
to:  
\begin{align*}
  \beta_0 \sim N(0, C), \beta_1 \sim N(0, C), \beta_2, \dots, \beta_m
  \overset{\text{i.i.d}}{\sim} N(0, \tau^2) ~~ \text{ and }
  \log \sigma \sim \text{uniform}(-C, C). 
\end{align*}
Therefore, we are bringing in a new parameter $\tau$ which controls
the scale of $\beta_2, \dots, \beta_m$. For now, treat $\tau$ as
fixed. So the prior on $\beta$ and $\sigma$ is now: 
\begin{align*}
  f_{\beta, \sigma}(\beta, \sigma) \propto  \frac{1}{\sigma}
  \frac{1}{\sqrt{\det Q}} \exp \left(-\frac{1}{2} \beta^T Q^{-1} \beta \right)
\end{align*}
where $Q$ is the $(m+1) \times (m+1)$ diagonal matrix with diagonal
entries $C, C, \tau^2, \dots, \tau^2$. 

The likelihood is unchanged: 
\begin{equation*}
  \left(\frac{1}{\sqrt{2 \pi}} \right)^n \sigma^{-n} \exp
  \left(-\frac{1}{2 \sigma^2} \|y - X \beta\|^2 \right). 
\end{equation*}
So the posterior for $\beta, \sigma$ is: 
\begin{align*}
  f_{\beta,\sigma \mid \text{data}}(\beta, \sigma)
 & \propto \frac{\sigma^{-n-1}}{\sqrt{\det Q}} \exp
   \left(-\frac{1}{2} \left(\frac{1}{\sigma^2} \|y - X \beta\|^2 +
   \beta^T Q^{-1} \beta \right) \right). 
\end{align*}
Using the completing the square formula: 
\begin{align*}
  \frac{1}{\sigma^2} \|y - X \beta\|^2 +
   \beta^T Q^{-1} \beta &= \left( \beta - \hat{\beta}\right)^T
                          \left(\frac{X^T X}{\sigma^2} + Q^{-1}
                          \right) \left(\beta - \hat{\beta} \right) \\
  &+ \frac{y^T y}{\sigma^2} - \left(\frac{y^T X}{\sigma^2} \right)
    \left(\frac{X^T X}{\sigma^2} + Q^{-1} \right)^{-1} \left(\frac{X^T
    y}{\sigma^2} \right), 
\end{align*}
where
\begin{align*}
  \hat{\beta} = \left(\frac{X^T X}{\sigma^2} + Q^{-1} \right)^{-1}
  \frac{X^T y}{\sigma^2}
\end{align*}
one can then show that
\begin{align}\label{postQ}
  \beta \mid \text{data}, \sigma, \tau \sim N\left( \left(\frac{X^T
                          X}{\sigma^2} + Q^{-1}  \right)^{-1} \frac{X^T y}{\sigma^2}, \left(\frac{X^T 
  X}{\sigma^2} + 
   Q^{-1} \right)^{-1} \right)  
\end{align}
Note that when $Q \rightarrow \infty$ (i.e., when $C \rightarrow
\infty$ and $\tau \rightarrow \infty$), this reverts to$N((X^T
X)^{-1} X^T y, \sigma^2 (X^T X)^{-1})$. To get the posterior of
$\sigma$ given $\tau$, we simply integrate $\beta$ from the joint
posterior to obtain: 

Integrating $\beta$ from the joint posterior gives the posterior of
$\gamma, \sigma$:
\begin{align*}
&  f_{\sigma \mid \text{data}, \tau}(\sigma) \\ &\propto
  \frac{\sigma^{-n-1}}{\sqrt{\det Q}}  \sqrt{\det
  \left(\frac{X^T X}{\sigma^2} + Q^{-1} \right)^{-1}} \exp \left(-\frac{y^T y}{2 \sigma^2} \right)\exp
  \left(\frac{y^T X}{2\sigma^2} \left(\frac{X^T 
  X}{\sigma^2} + 
   Q^{-1} \right)^{-1} \frac{X^T y}{\sigma^2}  \right). 
\end{align*}
Simplifying by using $\det Q \propto (\tau^2)^{m-1}$ and
$Q^{-1} \approx J/\tau^2$ where $J$ is the $(m+1) \times
(m+1)$ diagonal matrix with diagonals $0, 0, 1, \dots, 1$, we get
\begin{align*}
  &  f_{\sigma \mid \text{data}, \tau}(\sigma) \\ &\propto
\frac{\sigma^{-n+m}}{\tau^{m-1}} \left|X^T X + \frac{\sigma^2}{\tau^2}
                                                    J 
  \right|^{-1/2} \exp 
  \left(-\frac{y^T y - y^T X \left(X^T X + 
 \sigma^2\tau^{-2} J \right)^{-1} X^T y}{2\sigma^2} 
  \right) 
\end{align*}
This distribution is not easy to handle because of the presence of the
term $\sigma^2/\tau^2$ inside the determinant as well as inside the
inverse (in the exponent). One trick to simplify this is to
reparametrize and assume $\tau = \sigma \gamma$. With this, the above
conditional density becomes
\begin{align*}
  &  f_{\sigma \mid \text{data}, \gamma}(\sigma) \\ &\propto
\frac{1}{\gamma^{m-1}\sigma^{n-1}} \left|X^T X + \gamma^{-2} J 
  \right|^{-1/2} \exp 
  \left(-\frac{y^T y - y^T X \left(X^T X + 
\gamma^{-2} J \right)^{-1} X^T y}{2\sigma^2} 
  \right) 
\end{align*}
Ignoring terms above which do not depend on $\gamma$, we get 
\begin{equation*}
  f_{\sigma \mid \text{data}, \gamma}(\sigma) \propto \sigma^{-n+1} \exp 
  \left(-\frac{y^T y - y^T X \left(X^T X + 
 \gamma^{-2} J \right)^{-1} X^T y}{2\sigma^2} 
  \right). 
\end{equation*}
The right hand side above is actually the pdf of an inverse gamma
density (see
\url{https://en.wikipedia.org/wiki/Inverse-gamma_distribution}). This
can be seen by converting it into the density of $1/\sigma^2$:
\begin{align*}
  f_{1/\sigma^2 \mid \text{data}, \gamma}(x) &\propto f_{\sigma \mid
  \text{data}, \gamma}(x^{-1/2}) x^{-3/2} \\
  &\propto \left(x^{-1/2} \right)^{-n+1} \exp 
  \left(-x \frac{y^T y - y^T X \left(X^T X + 
 \gamma^{-2} J \right)^{-1} X^T y}{2} 
    \right) x^{-3/2} \\
  &= x^{(n-4)/2} \exp 
  \left(-x \frac{y^T y - y^T X \left(X^T X + 
 \gamma^{-2} J \right)^{-1} X^T y}{2} 
    \right). 
\end{align*}
Thus
\begin{align}\label{sig_post}
  \frac{1}{\sigma^2} \mid \text{data}, \gamma \sim
  \text{Gamma}\left(\frac{n}{2} - 1, \frac{y^T y - y^T X \left(X^T X + 
 \gamma^{-2} J \right)^{-1} X^T y}{2}. 
  \right) 
\end{align}
Thus when $\gamma$ is fixed, we can perform inference on $\beta,
\sigma$ using the above closed form formulae. Since $\gamma$ is also
unknown, we can place a prior $p(\gamma)$ on it, and then calculate
its posterior.  

Finally, we can marginalize $\sigma$ to obtain the posterior of
$\gamma$ alone as follows:
\begin{align*}
  f_{\gamma \mid \text{data}}(\gamma) &\propto \int_0^{\infty}
                                        p(\gamma)
                                        \frac{1}{\gamma^{m-1} \sigma^{n-1}}
 \left|X^T X + \gamma^{-2} J 
  \right|^{-1/2} \exp 
  \left(-\frac{y^T y - y^T X \left(X^T X + 
 \gamma^{-2} J \right)^{-1} X^T y}{2\sigma^2} 
                                        \right)  d\sigma  \\
&= p(\gamma)  \gamma^{-m+1} \left|X^T X + \gamma^{-2} J 
  \right|^{-1/2} \int_0^{\infty}  \sigma^{-n+1}\exp 
  \left(-\frac{y^T y - y^T X \left(X^T X + 
 \gamma^{-2} J \right)^{-1} X^T y}{2\sigma^2} 
                                        \right)  d\sigma. 
\end{align*}
Letting
\begin{align*}
  A := y^T y - y^T X \left(X^T X + 
 \gamma^{-2} J \right)^{-1} X^T y, 
\end{align*}
we get
\begin{align*}
  f_{\gamma \mid \text{data}}(\gamma) &\propto  p(\gamma) \gamma^{-m+1} \left|X^T X + \gamma^{-2} J 
  \right|^{-1/2} \int_0^{\infty} \sigma^{-n+1} \exp \left(-\frac{A}{2
                                        \sigma^2} \right) d\sigma 
\end{align*}
By the change of variable $\sigma = s \sqrt{A}$, we obtain
\begin{align*}
  f_{\gamma \mid \text{data}}(\gamma) &\propto  p(\gamma)  \gamma^{-m+1} \left|X^T X + \gamma^{-2} J 
  \right|^{-1/2} A^{-(n/2) + 1} \int_0^{\infty} s^{-n+1} \exp \left(-\frac{1}{2s^2}
  \right) ds \\
  &\propto p(\gamma)  \gamma^{-m+1} \left|X^T X + \gamma^{-2} J 
    \right|^{-1/2} A^{-(n/2) + 1} \\
  &=   \frac{p(\gamma)\gamma^{-m+1} \left|X^T X + \gamma^{-2} J 
    \right|^{-1/2}}{\left(y^T y - y^T X \left(X^T X + 
 \gamma^{-2} J \right)^{-1} X^T y \right)^{(n/2) - 1}}. 
\end{align*}
We usually choose $p(\gamma)$ as:
\begin{align*}
p(\gamma) \propto \frac{1}{\gamma} I\{\text{low} < \gamma <
  \text{high}\}
\end{align*}
for two fixed values $\text{low}$ and $\text{high}$. These could be
the values of  $\gamma$ which lead to the posterior mean (which
coincides with ridge regression) leading to underfitting and
overfitting respectively. 

Inference can be carried out by first taking a grid of
$\gamma$ values and computing the above posterior (on the logarithmic
scale) at the grid points. This posterior can be used to obtain
posterior samples of $\gamma$. For each sample of $\gamma$, we can
then sample $\sigma$ using the distribution \eqref{sig_post}. Given
samples from both $\gamma$ and $\sigma$, we can then sample $\beta$
using \eqref{postQ}.

\section{Induced Prior for the Regression Function}

Our regression function is being modeled as:
\begin{align}\label{reg_func}
  f(x) = \beta_0 + \beta_1 x + \beta_2 (x - 1)_+ + \dots + \beta_m(x -
  (m-1))_+
\end{align}
where $\beta_0, \beta_1, \dots, \beta_m$ are all independent with
\begin{align*}
  \beta_0, \beta_1 \overset{\text{i.i.d}}{\sim} N(0, C) ~~ \text{ and
  } ~~ \beta_2, \dots, \beta_m \overset{\text{i.i.d}}{\sim} N(0,
  \tau^2). 
\end{align*}
This can also be seen as a prior on the regression function $f$ more
directly. For every $0 \leq u_1 < \dots < u_k < \infty$, the joint
distribution of $(f(u_1), \dots, f(u_k))$ induced by \eqref{reg_func}
will be multivariate normal. This implies that $(f(u), u \geq 0)$ is a
Gaussian Process. The description of the GP can be done in terms of
the mean function $\E f(u)$ and the covariance kernel
$\text{Cov}(f(u), f(v))$.

The mean function is clearly zero (because each $\E \beta_j = 0$). The
covariance kernel is given by:
\begin{align*}
  &\text{cov}(f(u), f(v)) \\ &= \text{cov} \left(\beta_0 + \beta_1u  + \beta_2 (u - 1)_+ + \dots + \beta_m(u -
  (m-1))_+, \beta_0 + \beta_1 v + \beta_2 (v - 1)_+ + \dots + \beta_m(v -
                               (m-1))_+ \right) \\
  &= C(1 + uv) + \tau^2 \sum_{j=1}^{m-1} (u - j)_+ (v - j)_+ \\
  &= C(1 + uv) + \tau^2 \sum_{j=1}^{k} (u -
    j)(v - j) \\
  &= C(1 + uv) + \tau^2 \left(u v k + \frac{1}{6} k(k+1)(2k+1) -
    (u+v)\frac{k(k+1)}{2} \right) 
\end{align*}
where $k = \lfloor u \wedge v \rfloor$ (and $u \wedge v := \min(u,
v)$). 

Suppose now we use the approximation
\begin{align*}
  k \approx u \wedge v ~~ \text{ and } ~~ k(k+1)(2k+1) \approx 2k^3
  \approx 2(u \wedge v)^3 ~~ \text{ and } ~~ k(k+1) \approx k^2
  \approx (u \wedge v)^2. 
\end{align*}
Then the covariance kernel becomes:
\begin{align*}
  \text{cov}(f(u), f(v)) \approx C(1 + uv) + \tau^2 \left(u v (u
  \wedge v) + \frac{(u \wedge v)^3}{3} - \frac{(u \wedge v)^2}{2} (u + v)
  \right). 
\end{align*}
It turns out that the right hand side above is precisely the
covariance kernel of:
\begin{align*}
  G_x = \beta_0 + \beta_1 x + \tau I_x
\end{align*}
where $\beta_0, \beta_1 \overset{\text{i.i.d}}{\sim} N(0, C)$ and
$I_x$ is \textbf{Integrated Brownian Motion} on $[0, \infty)$. 
We will revisit this in the next lecture. 



%\begin{thebibliography}{99}

%\bibitem[MacKay and Peto(1995)]{MacKay1995HierarchicalDirichlet}
%D.~J.~C. MacKay and L.~C. Peto,
%\newblock ``A hierarchical Dirichlet language model,''
%\newblock \emph{Natural Language Engineering}, vol.~1,
%no.~3, pp.~289--308, 1995.

%\bibitem[Rubin(1981)]{Rubin1981BayesianBootstrap}
%D.~B. Rubin,
%\newblock ``The Bayesian bootstrap,''
%\newblock \emph{The Annals of Statistics}, vol.~9,
%no.~1, pp.~130--134, 1981.
  
%\bibitem[Fisher(1934)]{Fisher1934}
%R.~A. Fisher,
%\newblock ``Probability, likelihood and quantity of information in the logic of
%uncertain inference,''
%\newblock \emph{Proceedings of the Royal Society of London. Series~A,
%Containing Papers of a Mathematical and Physical Character}, vol.~146,
%no.~856, pp.~1--8, 1934.

%\bibitem[Efron(2010)]{Efron}
%Efron, B. (2010)
%\textit{Large-scale inference: empirical Bayes methods for estimation,
%testing, and prediction}. Cambridge University Press, 2010.

%\bibitem[Shumway and Stoffer(2017)]{ShumwayStoffer}
%Shumway, R. H. and Stoffer, D. S. (2017)
%\textit{Time Series Analysis and Its Applications: With R
%  Examples}. Springer, 4th edition. 
  
%     \bibitem[Jaynes(2003)]{Jaynes} Jaynes, E. T, (2003)
%  \textit{Probability theory: the logic of science}. Cambridge
%  University Press, 2003.
 
%  \bibitem[Sinai(1992)]{Sinai} Sinai, Y. G, (1992)
%  \textit{Probability theory: an introductory course}. Springer
%  Verlag, 1992.  

%\bibitem[{deGroot(1986)}]{DeGrootBlackwell}DeGroot, M. H. (1986)
%       A conversation with David Blackwell.
%\newblock \emph{Statistical Science}, \textbf{1}, 40--53.

%         \bibitem[Mosteller(1987)]{Mosteller} Mosteller, F, (1987)
%  \textit{Fifty challenging problems in probability with
%    solutions}. Courier Corporation, 1987.

%    \bibitem[MacKay(2003)]{MacKay} MacKay, D. J. C., (2003)
%  \textit{Information theory, inference, and learning
%  algorithms}. Cambridge University Press, 2003.

%\bibitem[{Lindley(2013)}]{Lindley}Lindley, D. V. (2013)
%   \textit{Understanding Uncertainty}. John Wiley and sons, 2013.


%\end{thebibliography}







\end{document}

%%% Local Variables:
%%% mode: latex
%%% TeX-master: t
%%% End:
